% Wersja źródłowa: audyt II, jednostka 09
\setcounter{section}{8}
\section{Banach--Alaoglu, Krein--Milman i topologia Mackeya}\label{sec:alaoglu-mackey}

\subsection{Zwartość słaba-* i równociągłość}

W przestrzeni nieskończenie wymiarowej kula jednostkowa nie jest zwarta
w topologii normowej. Zmiana topologii na słabą-* pozwala odzyskać zwartość,
ale zasadniczą rolę odgrywają wtedy \emph{zbiory polarne otoczeń zera}.
Najpierw dowiedziemy twierdzenia w klasie dowolnych przestrzeni
liniowo-topologicznych, bez zakładania lokalnej wypukłości i zupełności.

W całym rozdziale $\K\in\{\R,\C\}$. Dla TVS $E$ symbol $E'$ oznacza
przestrzeń sprzężoną, a $E^*$ przestrzeń algebraicznie sprzężoną.
Zbiorem polarnym $A\subseteq E$ nazywamy
$A^\circ=\{f\in E':|f(x)|\le1\ \text{dla każdego }x\in A\}$.
Na $E'$ rozpatrujemy topologię $\sigma(E',E)$ punktowej zbieżności na $E$.

\begin{definicja}[Równociągła rodzina funkcjonałów]\label{def:equicont-family}
Niech $E$ będzie TVS i $H\subseteq E'$. Rodzinę $H$ nazywamy
\emph{równociągłą}, jeżeli dla każdego $\varepsilon>0$ istnieje otoczenie
$U$ zera w $E$ takie, że $|f(u)|<\varepsilon$ dla wszystkich
$f\in H$ i $u\in U$.
\end{definicja}

\begin{stwierdzenie}[Równociągłość i zbiory polarne]\label{prop:equicont-polar}
Niech $E$ będzie TVS, a $H\subseteq E'$. Rodzina $H$ jest równociągła
wtedy i tylko wtedy, gdy istnieje otoczenie $U$ zera w $E$ takie,
że $H\subseteq U^\circ$. Każda rodzina równociągła jest punktowo
ograniczona, tzn. $\sup_{f\in H}|f(x)|<\infty$ dla każdego $x\in E$.
\end{stwierdzenie}
\begin{proof}
Jeżeli $H$ jest równociągła, bierzemy w definicji $\varepsilon=1$;
odpowiednie $U$ spełnia $H\subseteq U^\circ$. Odwrotnie,
jeżeli $H\subseteq U^\circ$, to dla ustalonego $\varepsilon>0$
otoczenie $(\varepsilon/2)U$ daje
$|f(v)|\le\varepsilon/2<\varepsilon$ dla $v\in(\varepsilon/2)U$
i $f\in H$. Dla $x\in E$ pochłanianie przez $U$ zapewnia
$x\in tU$ dla pewnego $t>0$, zatem $|f(x)|\le t$ dla $f\in H$.
\end{proof}

\begin{uwaga}[Wykorzystywane twierdzenie topologiczne]
W dowodzie twierdzenia Banacha--Alaoglu zastosujemy \emph{twierdzenie
Tychonowa}: iloczyn dowolnej rodziny zwartych przestrzeni topologicznych,
wyposażony w topologię produktową, jest zwarty. Jest to klasyczne
wynikanie z aksjomatu wyboru; nie będziemy tu ponownie dowodzić
tego ogólnotopologicznego rezultatu. Stosujemy go jedynie do iloczynów
zwartych przedziałów lub domkniętych kół w $\K$.
\end{uwaga}

\begin{twierdzenie}[Banach--Alaoglu--Bourbaki]\label{thm:banach-alaoglu-general}
Niech $E$ będzie przestrzenią liniowo-topologiczną nad $\K$, a
$U\subseteq E$ otoczeniem zera. Wówczas zbiór polarny $U^\circ$
jest zwarty w topologii $\sigma(E',E)$.
Nie zakładamy, że $E$ jest Hausdorffa, lokalnie wypukła, metryzowalna
lub zupełna.
\end{twierdzenie}
\begin{proof}
Umieszczamy $E^*$ w $\K^E$, utożsamiając funkcjonał z rodziną
jego wartości $(f(x))_{x\in E}$. Topologia punktowej zbieżności
jest topologią indukowaną przez topologię produktową.
Ponieważ $U$ jest pochłaniające, dla każdego $x\in E$ można wybrać
$t_x>0$ takie, że $x\in t_xU$. Każdy $f\in U^\circ$ spełnia wtedy
$|f(x)|\le t_x$. Niech $D_x=\{z\in\K:|z|\le t_x\}$ i
$P=\prod_{x\in E}D_x$. Zbiory $D_x$ są zwarte, więc $P$ jest zwarty
na mocy twierdzenia Tychonowa.

Zbiór $U^\circ$ jest domknięty w $P$. Rzeczywiście, warunki
$f(x+y)=f(x)+f(y)$, $f(ax)=af(x)$ oraz $|f(u)|\le1$ dla $u\in U$
opisują przecięcia zbiorów domkniętych w topologii produktowej.
Każda funkcja $f:E\to\K$ spełniająca te warunki jest liniowa,
a przy tym ciągła: dla $\varepsilon>0$ na otoczeniu
$(\varepsilon/2)U$ zachodzi $|f(v)|\le\varepsilon/2$.
Zatem otrzymany zbiór jest dokładnie $U^\circ\subseteq E'$.
Jako domknięty podzbiór zwartego $P$ jest on zwarty. Topologia
indukowana z $P$ na $U^\circ$ to $\sigma(E',E)$.
\end{proof}

\begin{wniosek}[Kula przestrzeni sprzężonej]\label{cor:alaoglu-norm}
Niech $E$ będzie przestrzenią unormowaną, niekoniecznie zupełną.
Dla każdego $R\ge0$ kula
$B_R'=\{f\in E':\|f\|\le R\}$ jest zwarta w $\sigma(E',E)$.
W szczególności normowa kula jednostkowa przestrzeni sprzężonej
każdej przestrzeni Banacha jest słabo-* zwarta.
\end{wniosek}
\begin{proof}
Dla $R>0$ niech $U=\{x\in E:\|x\|<1/R\}$.
Wtedy $U^\circ=B_R'$: warunek $|f(u)|\le1$ dla $\|u\|<1/R$
jest równoważny $\|f\|\le R$. Zastosowanie
Twierdzenia~\ref{thm:banach-alaoglu-general} kończy dowód.
Dla $R=0$ zbiór $B_0'=\{0\}$ jest zwarty.
\end{proof}

\begin{wniosek}[Domknięte podzbiory kuli]\label{cor:alaoglu-closed}
Niech $E$ będzie przestrzenią unormowaną, $R\ge0$, a
$A\subseteq B_R'$ będzie domknięty w topologii
$\sigma(E',E)$ indukowanej na $B_R'$. Wtedy $A$ jest
słabo-* zwarty.
\end{wniosek}
\begin{proof}
Domknięty podzbiór zwartej przestrzeni jest zwarty.
\end{proof}

\begin{uwaga}[Czego twierdzenie nie mówi?]
Twierdzenie Banacha--Alaoglu nie orzeka, że cała przestrzeń $E'$
jest słabo-* zwarta ani że kula jest zwarta w normie.
Zwartość nie daje też automatycznie możliwości zastąpienia
sieci przez ciągi: do tego potrzebna jest dodatkowa metryzowalność.
\end{uwaga}

\subsection{Metryzowalność i kryterium ośrodkowości}

\begin{twierdzenie}[Metryzowalność na rodzinie równociągłej]
\label{thm:weakstar-metrizable-equicont}
Niech $E$ będzie ośrodkową TVS Hausdorffa, a $H\subseteq E'$
rodziną równociągłą. Wówczas topologia indukowana na $H$
przez $\sigma(E',E)$ jest metryzowalna.
\end{twierdzenie}
\begin{proof}
Niech $D=\{x_n:n\in\N\}$ będzie zbiorem gęstym w $E$.
Określamy na $H$ funkcję
\[
 d(f,g)=\sum_{n=1}^\infty 2^{-n}\min\{1,|f(x_n)-g(x_n)|\}.
\]
Jest to metryka: nierówność trójkąta wynika z
$\min\{1,a+b\}\le\min\{1,a\}+\min\{1,b\}$,
a jeśli $d(f,g)=0$, to $f$ i $g$ pokrywają się na gęstym
$D$, zatem przez ciągłość na całym $E$.
Topologia metryki jest topologią punktowej zbieżności na $D$,
więc jest nie silniejsza od $\sigma(E',E)|_H$.

Pokażemy odwrotną inkluzję. Ustalmy $x\in E$, $f_0\in H$
i $\varepsilon>0$. Ponieważ $H$ jest równociągła,
istnieje zrównoważone otoczenie $V$ zera takie, że
$|h(v)|<\varepsilon/3$ dla wszystkich $h\in H$, $v\in V$.
Wybieramy $x_n\in D\cap(x+V)$. Dla $f\in H$ mamy
\[
 |(f-f_0)(x)|\le |f(x-x_n)|+|(f-f_0)(x_n)|
                      +|f_0(x_n-x)|
 <2\varepsilon/3+|(f-f_0)(x_n)|.
\]
Warunek $|(f-f_0)(x_n)|<\varepsilon/3$ jest otwarty
w metryce $d$. Zatem wartościowania w każdym $x\in E$
są ciągłe na $(H,d)$, co daje odwrotną inkluzję topologii.
\end{proof}

\begin{twierdzenie}[Kryterium dla kuli przestrzeni sprzężonej]
\label{thm:banach-ball-metrizable-iff-separable}
Niech $E$ będzie przestrzenią unormowaną nad $\K$,
a $B'=\{f\in E':\|f\|\le1\}$. Wówczas następujące
warunki są równoważne:
\begin{enumerate}
\item $E$ jest ośrodkowa w topologii normowej;
\item $B'$ z topologią $\sigma(E',E)$ jest metryzowalna.
\end{enumerate}
Jeśli warunki zachodzą, to $B'$ jest zwarta, a każdy jej ciąg ma podciąg zbieżny.
\end{twierdzenie}
\begin{proof}
Kula $B'$ jest równociągła, bo $|f(x)|\le\|x\|$ dla $f\in B'$.
Jeżeli $E$ jest ośrodkowa, stosujemy
Twierdzenie~\ref{thm:weakstar-metrizable-equicont};
zwartość wynika z Twierdzenia~\ref{thm:banach-alaoglu-general}.
Każda zwarta przestrzeń metryczna ma własność zwartości ciągowej.

Odwrotnie, załóżmy, że $B'$ jest metryzowalna. Istnieje
przeliczalna baza otoczeń zera w $B'$, a więc można wybrać
ciąg podstawowych otoczeń słabych-* zera
$V_n=\{f\in B':|f(x)|<\varepsilon_n\ \text{dla }x\in F_n\}$,
gdzie $F_n\subseteq E$ jest skończony, tak aby $\{V_n\}$
była bazą otoczeń zera w $B'$. Niech
$M=\operatorname{lin}(\bigcup_{n=1}^\infty F_n)$.
Jeśli $\overline M^{\|\cdot\|}\ne E$, twierdzenie Hahna--Banacha
daje niezerowy $f\in B'$ zerujący się na $M$.
Wtedy $f\in V_n$ dla każdego $n$. Ponieważ $B'$ jest Hausdorffa,
przecięcie wszystkich otoczeń zera w $B'$ jest równe $\{0\}$;
sprzeczność. Zatem $M$ jest gęsta w $E$, a jako przestrzeń
rozpięta nad $\K$ przez zbiór przeliczalny ma przeliczalny
gęsty podzbiór (kombinacje wymierne, a nad $\C$ -- z
wymiernymi częściami rzeczywistymi i urojonymi współczynników).
Stąd $E$ jest ośrodkowa.
\end{proof}

\begin{przyklad}[Kula przestrzeni dualnej $\ell^1$]
Dla $E=c_0$ mamy $E'=\ell^1$. Ponieważ $c_0$ jest ośrodkowa,
kula jednostkowa w $\ell^1$ jest zwarta i metryzowalna
w $\sigma(\ell^1,c_0)$. Nie jest natomiast zwarta w normie:
wektory $e_n$ należą do kuli, lecz $\|e_n-e_m\|_1=2$
dla $n\ne m$. Przy tym $e_n\to0$ słabo-*.
\end{przyklad}

\subsection{Dlaczego sieci są niezbędne?}

\begin{przyklad}[Zwartość bez metryzowalności]\label{ex:uncountable-product}
Niech $I$ będzie zbiorem nieprzeliczalnym,
$E=\ell^1(I)$, a $E'=\ell^\infty(I)$ z kanonicznym parowaniem
$\langle x,f\rangle=\sum_{i\in I}x_i f_i$.
Kula jednostkowa $B'$ przestrzeni $\ell^\infty(I)$
wyposażona w $\sigma(\ell^\infty(I),\ell^1(I))$
jest homeomorficzna z produktem
$\prod_{i\in I}\{z\in\K:|z|\le1\}$.
Istotnie, zbieżność słaba-* implikuje współrzędną przez
podstawienie $x=e_i$. Odwrotnie, jeśli sieć $f_\alpha\in B'$
zbiega współrzędnie do $f$, to dla $x\in\ell^1(I)$
wybieramy skończony $F\subseteq I$ z
$\sum_{i\notin F}|x_i|<\varepsilon/4$ i kontrolujemy
skończoną sumę $\sum_{i\in F}x_i(f_{\alpha,i}-f_i)$;
reszta jest mniejsza od $2\sum_{i\notin F}|x_i|<\varepsilon/2$.
Produkt jest zwarty, lecz nie jest metryzowalny, ponieważ
$\ell^1(I)$ nie jest ośrodkowa.

Wewnątrz $B'$ rozpatrzmy zbiór
$A=\{\chi_F:F\subseteq I\ \text{jest skończony}\}$,
gdzie $\chi_F$ jest funkcją charakterystyczną $F$.
Sieć $(\chi_F)$, indeksowana skończonymi podzbiorami $I$
i uporządkowana inkluzją, zbiega słabo-* do funkcji
$\chi_I$. Żaden \emph{ciąg} elementów $A$ nie zbiega
jednak słabo-* do $\chi_I$: suma jego przeliczalnie wielu
skończonych nośników jest przeliczalna, a więc istnieje
$i\in I$, na którym wszystkie wyrazy ciągu mają wartość zero.
\end{przyklad}

\begin{przyklad}[Słabo-* zwarta rodzina nierównociągła]
\label{ex:compact-not-equicont}
Niech $E=c_{00}$ będzie przestrzenią ciągów o skończonym nośniku,
z normą $\|x\|_\infty=\sup_n|x_n|$. Przestrzeń $E$ nie jest
zupełna. Dla $n\in\N$ określamy $f_n\in E'$ wzorem
$f_n(x)=n x_n$. Ponieważ każdy $x\in c_{00}$ ma skończony
nośnik, $f_n(x)=0$ dla wszystkich dostatecznie dużych $n$.
Wobec tego $f_n\to0$ w $\sigma(E',E)$, a zbiór
$K=\{0\}\cup\{f_n:n\in\N\}$ jest słabo-* zwarty:
każdy otwarty zbiór zawierający zero zawiera prawie wszystkie $f_n$.
Jednak $\|f_n\|=n$, zatem $K$ nie jest równociągły.
To pokazuje, że słaba-* zwartość nie oznacza równociągłości
bez dodatkowych założeń. W rozdziale o twierdzeniu
Banacha--Steinhausa omówimy warunki, pod którymi implikacja
staje się prawdziwa.
\end{przyklad}

\subsection{Punkty ekstremalne i twierdzenie Kreina--Milmana}
\label{subsec:krein-milman}

Twierdzenie Banacha--Alaoglu dostarcza zwartych zbiorów wypukłych
w topologiach słabych-*. Kolejne twierdzenie opisuje ich geometrię:
cały zbiór można odzyskać z jego punktów ekstremalnych. W tej
części, jeśli nie zaznaczono inaczej, $E$ oznacza przestrzeń
lokalnie wypukłą Hausdorffa nad $\K\in\{\R,\C\}$. W przypadku
zespolonym wypukłość rozumiemy w sensie rzeczywistym.

\begin{definicja}[Punkt ekstremalny i ściana]
\label{def:extreme-face9}
Niech $K\subseteq E$ będzie niepustym zbiorem wypukłym.
Punkt $e\in K$ jest \emph{ekstremalny}, jeżeli z
$e=tx+(1-t)y$, $x,y\in K$, $0<t<1$ wynika $x=y=e$.
Zbiór wszystkich takich punktów oznaczamy $\operatorname{Ext}(K)$.
Niepusty podzbiór $F\subseteq K$ jest \emph{ścianą} zbioru $K$,
jeżeli jest wypukły i dla $x,y\in K$ oraz $0<t<1$ z
$tx+(1-t)y\in F$ wynika $x,y\in F$.
Przez $\operatorname{conv}(A)$ oznaczamy otoczkę wypukłą $A$,
a przez $\overline A$ jej domknięcie w $E$.
\end{definicja}

\begin{przyklad}[Wielokąt i kula w przestrzeni Hilberta]
Wierzchołki wypukłego wielokąta o niepustym wnętrzu w $\R^2$ są dokładnie jego
punktami ekstremalnymi. Jeżeli $H$ jest rzeczywistą przestrzenią
Hilberta, to punktami ekstremalnymi domkniętej kuli jednostkowej
$B_H$ są dokładnie punkty sfery $\{x:\|x\|=1\}$.
Istotnie, dla $\|x\|<1$ i dostatecznie małego $h\ne0$
mamy $x=\tfrac12(x+h)+\tfrac12(x-h)$, a oba punkty należą do kuli.
Natomiast jeśli $\|x\|=1$ i $x=(u+v)/2$ przy $u,v\in B_H$,
to tożsamość równoległoboku daje
$\|u-v\|^2=2\|u\|^2+2\|v\|^2-4\|x\|^2\le0$, więc $u=v=x$.
\end{przyklad}

\begin{center}
\begin{tikzpicture}[scale=0.90, every node/.style={font=\small}]
\filldraw[fill=blue!7,draw=blue!55!black,very thick]
(0.1,0.6)--(1.5,0.15)--(3.9,0.65)--(4.35,2.05)--(2.75,2.85)--(0.4,2.3)--cycle;
\foreach \x/\y in {0.1/0.6,1.5/0.15,3.9/0.65,4.35/2.05,2.75/2.85,0.4/2.3}
 {\fill[blue!65!black] (\x,\y) circle (2.7pt);}
\draw[orange!85!black,densely dashed,thick] (-0.1,2.85)--(4.75,2.85);
\fill[orange!80!black] (2.75,2.85) circle (3.5pt);
\node[anchor=west,align=left] at (4.83,2.9) {$\ell=\max_K\ell$};
\node at (2.2,1.5) {$K$};
\node[anchor=north] at (2.2,-0.05) {Wierzchołki są punktami ekstremalnymi.};
\end{tikzpicture}
\end{center}
\noindent
\emph{Ilustracja.} Linia przerywana jest poziomicą funkcjonału
liniowego podpierającą $K$; odpowiednia ściana składa się tutaj
z jednego punktu ekstremalnego.

\begin{lemat}[Ściana poziomicy maksymalnej]
\label{lem:max-face9}
Niech $K\subseteq E$ będzie niepusty, zwarty i wypukły,
a $\ell:E\to\R$ niech będzie ciągłym rzeczywistym
funkcjonałem liniowym. Wówczas zbiór
$F_\ell=\{x\in K:\ell(x)=\max_{y\in K}\ell(y)\}$
jest niepustą zwartą ścianą $K$.
Jeżeli $F$ jest ścianą $K$, a $G$ ścianą $F$,
to $G$ jest również ścianą $K$.
\end{lemat}
\begin{proof}
Ciągłość $\ell$ i zwartość $K$ zapewniają istnienie maksimum;
$F_\ell$ jest domkniętym, a więc zwartym podzbiorem $K$.
Jest wypukły z liniowości $\ell$. Jeśli
$tx+(1-t)y\in F_\ell$, gdzie $x,y\in K$ i $0<t<1$, to
$\max_K\ell=t\ell(x)+(1-t)\ell(y)$, podczas gdy
$\ell(x),\ell(y)\le\max_K\ell$. Stąd obie wartości
są maksymalne i $x,y\in F_\ell$. Druga część wynika
z dwukrotnego zastosowania definicji ściany.
\end{proof}

\begin{lemat}[Istnienie punktu ekstremalnego]
\label{lem:exist-extreme9}
Każdy niepusty, zwarty i wypukły podzbiór przestrzeni lokalnie
wypukłej Hausdorffa ma co najmniej jeden punkt ekstremalny.
Co więcej, każda jego niepusta, domknięta ściana zawiera
punkt ekstremalny całego zbioru.
\end{lemat}
\begin{proof}
Niech $K\subseteq E$ spełnia założenia. Rozważmy rodzinę
$\mathcal F$ wszystkich niepustych \emph{domkniętych} ścian $K$,
uporządkowaną przez inkluzję odwrotną.
Każdy łańcuch w $\mathcal F$ ma górne ograniczenie:
jego elementy są zagnieżdżonymi niepustymi zbiorami
domkniętymi w zwartym $K$, więc mają niepuste przecięcie.
Przecięcie jest domkniętą ścianą, ponieważ warunek definiujący
ścianę jest zachowywany przy przecięciach.
Lemat Kuratowskiego--Zorna zapewnia istnienie minimalnej
niepustej domkniętej ściany $F$.

Gdyby $F$ zawierała różne punkty $x,y$, lokalna wypukłość
i warunek Hausdorffa zapewniałyby ciągły rzeczywisty
funkcjonał liniowy $\ell$ rozróżniający $x$ i $y$.
Dla $\K=\C$ możemy wziąć część rzeczywistą ciągłego
zespolonego funkcjonału. Z Lematu~\ref{lem:max-face9}
wynika, że zbiór punktów, w których $\ell$ osiąga
maksimum na $F$, jest niepustą, domkniętą ścianą $K$,
prawdziwie zawartą w $F$. To przeczy minimalności.
Zatem $F=\{e\}$, a warunek bycia ścianą oznacza,
że $e\in\operatorname{Ext}(K)$.
Jeśli zaczniemy od dowolnej niepustej domkniętej ściany
$F_0$ zamiast od $K$, ten sam argument da jej punkt
ekstremalny, ekstremalny również w $K$.
\end{proof}

\begin{twierdzenie}[Krein--Milman]
\label{thm:krein-milman}
Niech $E$ będzie rzeczywistą lub zespoloną przestrzenią
lokalnie wypukłą Hausdorffa, a $K\subseteq E$ niepustym,
zwartym i wypukłym zbiorem. Wtedy
\[
 K=\overline{\operatorname{conv}(\operatorname{Ext}(K))}.
\]
Domknięcie rozumiemy w wyjściowej topologii $E$.
\end{twierdzenie}
\begin{proof}
Z Lematu~\ref{lem:exist-extreme9} zbiór
$\operatorname{Ext}(K)$ jest niepusty. Ponieważ $K$ jest
wypukły i domknięty, zbiór
$C=\overline{\operatorname{conv}(\operatorname{Ext}(K))}$
jest zawarty w $K$ i jest zwarty jako domknięty podzbiór $K$.

Załóżmy, że istnieje $x_0\in K\setminus C$.
Twierdzenie o ścisłym oddzielaniu punktu od domkniętego
zbioru wypukłego w przestrzeni lokalnie wypukłej
(twierdzenie Hahna--Banacha o oddzielaniu zbiorów wypukłych) daje ciągły
rzeczywisty funkcjonał liniowy $\ell:E\to\R$ spełniający
$\ell(x_0)>\sup_{c\in C}\ell(c)$. W przypadku zespolonym
$\ell$ można otrzymać jako część rzeczywistą funkcjonału
zespolonego.

Niech $M=\max_{x\in K}\ell(x)$.
Z Lematu~\ref{lem:max-face9} zbiór
$F_\ell=\{x\in K:\ell(x)=M\}$ jest niepustą zwartą
ścianą $K$. Z Lematu~\ref{lem:exist-extreme9}
ma on punkt $e\in\operatorname{Ext}(K)$.
Wtedy $e\in C$, a równocześnie
$\ell(e)=M\ge\ell(x_0)>\sup_C\ell$, co jest sprzecznością.
Stąd $C=K$.
\end{proof}

\begin{uwaga}[Dlaczego hipotezy są istotne?]
Zwartość nie może być pominięta: cała przestrzeń $E=\R$
ma niepusty, domknięty i wypukły zbiór $K=\R$,
ale nie ma punktów ekstremalnych. Lokalna wypukłość
gwarantuje istnienie dostatecznie wielu ciągłych
funkcjonałów do oddzielania punktów. W ogólnych TVS
bez lokalnej wypukłości powyższy dowód nie działa.
Twierdzenie nie mówi także, że każdy punkt $K$ jest
\emph{skończoną} kombinacją wypukłą punktów ekstremalnych;
w ogólności konieczne jest domknięcie.
\end{uwaga}

\begin{wniosek}[Zasada maksimum Bauera]
\label{cor:bauer9}
Niech $K$ będzie niepustym zwartym wypukłym podzbiorem
przestrzeni lokalnie wypukłej Hausdorffa $E$ nad $\K$.
Jeżeli $f:K\to\R$ jest wypukła i półciągła z góry,
to istnieje $e\in\operatorname{Ext}(K)$ takie,
że $f(e)=\max_{x\in K}f(x)$.
\end{wniosek}
\begin{proof}
Niech $M=\max_K f$; maksimum istnieje dzięki zwartości $K$
i półciągłości $f$ z góry. Zbiór $A=\{x\in K:f(x)=M\}$
jest niepusty i domknięty. Nie musi być wypukły, ale ma
\emph{własność ekstremalności}: jeśli $tx+(1-t)y\in A$
przy $x,y\in K$ i $0<t<1$, to $x,y\in A$. Rzeczywiście,
$M=f(tx+(1-t)y)\le tf(x)+(1-t)f(y)\le M$,
więc obie nierówności są równościami.

Wybierzmy, korzystając z lematu Kuratowskiego--Zorna,
minimalny niepusty, domknięty podzbiór $B\subseteq A$
mający własność ekstremalności względem $K$.
Argument z przecięciem łańcucha jest taki sam jak
w Lemacie~\ref{lem:exist-extreme9}; nie używa wypukłości $B$.
Gdyby $B$ zawierał różne punkty, istniałby ciągły rzeczywisty
funkcjonał $\ell$ niebędący na nim stałym. Niepusty zbiór
$B_\ell=\{x\in B:\ell(x)=\max_{y\in B}\ell(y)\}$
byłby zwartym, właściwym podzbiorem $B$.
Jeśli $z=tx+(1-t)y\in B_\ell$ dla $x,y\in K$, to
$x,y\in B$, gdyż $B$ jest ekstremalny, a z liniowości
$\ell$ wynika $x,y\in B_\ell$. Jest to sprzeczność
z minimalnością. Zatem $B=\{e\}$; własność ekstremalności
tej jednoelementowej rodziny mówi dokładnie,
że $e\in\operatorname{Ext}(K)$, a $e\in A$ kończy dowód.
\end{proof}

\begin{przyklad}[Produkt przedziałów]
\label{ex:krein-cube9}
Niech $I$ będzie dowolnym niepustym zbiorem,
a $K=[-1,1]^I$ z topologią produktową, uważany za
podzbiór $\R^I$. Jest to zwarty zbiór wypukły na mocy
twierdzenia Tychonowa. Jego punkty ekstremalne to dokładnie
$x\in K$ takie, że $x_i\in\{-1,1\}$ dla każdego $i\in I$.
Istotnie, jeżeli $|x_i|<1$, można zmienić tylko współrzędną
$i$ o małe wartości $\pm\varepsilon$, aby przedstawić
$x$ jako środek dwóch różnych punktów $K$.
Jeśli wszystkie współrzędne są równe $\pm1$, to każdy
rozkład $x=ty+(1-t)z$ w $K$ musi być równy $y=z=x$
w każdej współrzędnej.
Dla dowolnego nieprzeliczalnego $I$ argument zwartości
opiera się na twierdzeniu Tychonowa, a nie na kryterium
sprowadzającym zwartość wyłącznie do zachowania ciągów.
\end{przyklad}

\begin{przyklad}[Słabo-* zwarta kula $\ell^1$]
\label{ex:krein-l1-weakstar9}
Na przestrzeni $\ell^1(\N;\R)=c_0(\N;\R)'$ rozważmy
słabą-* topologię $\sigma(\ell^1,c_0)$.
Kula $B=\{x\in\ell^1:\|x\|_1\le1\}$ jest zwarta
w tej topologii na mocy twierdzenia Banacha--Alaoglu.
Jej punktami ekstremalnymi są dokładnie $\pm e_n$.
Dla $\|x\|_1<1$ można rozdzielić $x$ jako średnią
$x\pm\varepsilon e_1$. Jeśli $\|x\|_1=1$ i $x$ ma
co najmniej dwie niezerowe współrzędne, odpowiednio małe
przeniesienie masy między nimi pozwala zapisać
$x=(y+z)/2$ dla różnych $y,z\in B$.
Z kolei $\pm e_n$ są ekstremalne, gdyż równość
$\pm1=(y_n+z_n)/2$ dla $y,z\in B$ wymusza
$y_n=z_n=\pm1$ i zerowanie się innych współrzędnych.
Wobec twierdzenia Kreina--Milmana kula $B$ jest słabo-*
domkniętą otoczką wypukłą tych punktów.
\end{przyklad}


\subsection{Twierdzenie Mackeya--Arensa}

Przypomnijmy, że dla pary dualnej $(E,F)$ topologia
$\tau(E,F)$ jest generowana przez seminormy
$p_K(x)=\sup_{f\in K}|\langle x,f\rangle|$, gdzie $K$
przebiega niepuste, absolutnie wypukłe i zwarte zbiory
w topologii $\sigma(F,E)$.

\begin{definicja}[Topologia zgodna z parą dualną]\label{def:compatible-duality}
Niech $(E,F,\langle\cdot,\cdot\rangle)$ będzie parą dualną.
Lokalnie wypukłą topologię Hausdorffa $\rho$ na $E$ nazywamy
\emph{zgodną z parą dualną $(E,F)$}, jeżeli odwzorowanie
$F\to(E,\rho)'$, $f\mapsto(x\mapsto\langle x,f\rangle)$,
jest izomorfizmem algebraicznym.
\end{definicja}

\begin{lemat}[Zwartość i dominacja funkcjonału]\label{lem:compact-domination}
Niech $(E,F)$ będzie parą dualną, a $K\subseteq F$ niepustym,
absolutnie wypukłym i zwartym zbiorem w $\sigma(F,E)$.
Jeżeli $g\in E^*$ spełnia
$|g(x)|\le p_K(x)$ dla każdego $x\in E$,
to istnieje $f\in K$ takie, że
$g(x)=\langle x,f\rangle$ dla wszystkich $x\in E$.
\end{lemat}
\begin{proof}
Dla każdego skończonego zbioru $S=\{x_1,\ldots,x_m\}\subseteq E$
rozważamy odwzorowanie
$T_S:K\to\K^m$, $T_S(f)=(\langle x_1,f\rangle,\ldots,
\langle x_m,f\rangle)$. Jego obraz jest zwarty i absolutnie wypukły.
Pokażemy, że $y=(g(x_1),\ldots,g(x_m))\in T_S(K)$.
Gdyby tak nie było, twierdzenie o silnym oddzielaniu punktu
od zwartego wypukłego zbioru w skończonym wymiarze dałoby
$a_1,\ldots,a_m\in\K$ takie, że
\[
 \operatorname{Re}\sum_{j=1}^m a_jg(x_j)
 >\sup_{f\in K}\operatorname{Re}
       \sum_{j=1}^m a_j\langle x_j,f\rangle.
\]
Połóżmy $x=\sum_{j=1}^m a_jx_j$.
Absolutna wypukłość $K$ oznacza, że zbiór wartości
$\{\langle x,f\rangle:f\in K\}$ jest zrównoważony,
a więc jego supremum części rzeczywistych jest równe
supremum modułów. Otrzymalibyśmy
$\operatorname{Re}g(x)>p_K(x)$, sprzecznie z założeniem.

Zatem dla każdego skończonego $S$ zbiór
$K_S=\{f\in K:\langle x,f\rangle=g(x)\ \text{dla }x\in S\}$
jest niepusty i domknięty w $K$. Rodzina $(K_S)$ ma
własność skończonych przecięć, ponieważ
$K_{S_1}\cap\cdots\cap K_{S_r}=K_{S_1\cup\cdots\cup S_r}$.
Zwartość $K$ daje punkt we wspólnym przecięciu,
co kończy dowód.
\end{proof}

\begin{twierdzenie}[Mackey--Arens]\label{thm:mackey-arens}
Niech $(E,F,\langle\cdot,\cdot\rangle)$ będzie parą dualną
przestrzeni liniowych nad $\K$. Wówczas topologia
$\tau(E,F)$ jest \emph{najsilniejszą} lokalnie wypukłą
topologią Hausdorffa na $E$ zgodną z tą parą dualną.
Dokładniej:
\begin{enumerate}
\item przestrzeń sprzężona $(E,\tau(E,F))'$ jest
      algebraicznie równa $F$ przez dane parowanie;
\item jeżeli $\rho$ jest dowolną inną lokalnie wypukłą topologią
      Hausdorffa na $E$ zgodną z $(E,F)$, to $\rho\subseteq\tau(E,F)$.
\end{enumerate}
\end{twierdzenie}
\begin{proof}
Z definicji $\tau(E,F)$ każdy element $F$ wyznacza
funkcjonał ciągły na $E$; istotnie, zbiór
$\{af:|a|\le1\}$ jest absolutnie wypukły i zwarty
w $\sigma(F,E)$.

Niech $g\in(E,\tau(E,F))'$. Ciągłość w zerze daje
niepuste absolutnie wypukłe zbiory zwarte
$K_1,\ldots,K_m\subseteq F$ i stałą $C>0$ takie, że
$|g(x)|\le C\max_j p_{K_j}(x)$ dla wszystkich $x\in E$.
Aby uzyskać to oszacowanie, bierzemy bazowe otoczenie
$\{x:\max_j p_{K_j}(x)<\varepsilon\}$ zawarte w
$\{x:|g(x)|<1\}$ i skalujemy $x$; można przyjąć
$C=2/\varepsilon$. Niech $K=K_1+\cdots+K_m$.
Jest to niepusty, absolutnie wypukły zbiór zwarty
w $\sigma(F,E)$ jako obraz zwartego produktu przez dodawanie.
Ponieważ $0\in K_j$, mamy $p_K\ge\max_j p_{K_j}$.
Zatem $|g(x)|\le C p_K(x)=p_{CK}(x)$.
Lemat~\ref{lem:compact-domination}, zastosowany do zwartego
absolutnie wypukłego $CK$, zapewnia $g\in F$.
To dowodzi punktu (1).

Niech teraz $\rho$ będzie zgodną lokalnie wypukłą topologią.
Wystarczy sprawdzić, że każda $\rho$-ciągła seminorma
$p:E\to[0,\infty)$ jest $\tau(E,F)$-ciągła.
Określmy $U=\{x:p(x)\le1\}$ i
\[
 K_p=\{f\in F:|\langle x,f\rangle|\le p(x)
                 \ \text{dla wszystkich }x\in E\}.
\]
Zbiór $U$ jest otoczeniem zera dla $\rho$, a
$K_p=U^\circ$ przy danym utożsamieniu $F=(E,\rho)'$.
Istotnie, jeśli $|\langle x,f\rangle|\le1$ dla $p(x)\le1$,
to dla $p(x)>0$ stosujemy $x/p(x)$, a przy $p(x)=0$
skalujemy $nx$ i przechodzimy z $n\to\infty$.
Na mocy Twierdzenia~\ref{thm:banach-alaoglu-general}
zbiór $K_p$ jest zwarty w $\sigma(F,E)$;
jest również niepusty i absolutnie wypukły.

Pokażemy, że $p=p_{K_p}$. Nierówność $p_{K_p}\le p$
wynika z definicji. Dla ustalonego $x\ne0$ funkcjonał
na $\K x$, $\lambda x\mapsto\lambda p(x)$,
jest dominowany przez seminormę $p$ co do modułu.
Twierdzenie Hahna--Banacha (w wersji dla seminormy)
przedłuża go do funkcjonału $f\in E^*$ z
$|f(y)|\le p(y)$ dla każdego $y\in E$.
Funkcjonał ten jest $\rho$-ciągły, należy więc do $F$,
a w konsekwencji $f\in K_p$ i $|f(x)|=p(x)$.
Zachodzi zatem $p=p_{K_p}$, czyli $p$ jest
$\tau(E,F)$-ciągła. Ponieważ topologia $\rho$ jest
generowana przez swoje ciągłe seminormy, otrzymujemy
$\rho\subseteq\tau(E,F)$.
\end{proof}

\begin{wniosek}[Najsłabsza i najsilniejsza topologia zgodna]\label{cor:mackey-extremes}
Niech $(E,F)$ będzie parą dualną. Topologia słaba
$\sigma(E,F)$ jest najsłabszą, zaś topologia Mackeya
$\tau(E,F)$ najsilniejszą lokalnie wypukłą topologią
Hausdorffa zgodną z tą parą. W szczególności dla każdej
takiej topologii $\rho$ zachodzi
$\sigma(E,F)\subseteq\rho\subseteq\tau(E,F)$.
\end{wniosek}
\begin{proof}
Jeżeli $\rho$ jest zgodna, wszystkie funkcjonały z $F$
są $\rho$-ciągłe, więc minimalność topologii początkowej
daje $\sigma(E,F)\subseteq\rho$. Druga inkluzja wynika
z Twierdzenia~\ref{thm:mackey-arens}.
\end{proof}

\begin{przyklad}[Przestrzeń słaba i topologia Mackeya nie muszą być równe]
Niech $E$ będzie nieskończenie wymiarową przestrzenią Banacha,
a $F=E'$. Topologia $\sigma(E,E')$ nie jest normowalna
(co udowodniono w \hyperref[sec:topologie-slabe]{rozdziale 8}), a topologia normowa
jest zgodna z parą $(E,E')$. Z wniosku
\ref{cor:mackey-extremes} wynika więc, że
$\sigma(E,E')\subsetneq\tau(E,E')$.
Nie utożsamiamy tu jednak automatycznie $\tau(E,E')$
z topologią normową: odpowiednie twierdzenie dla przestrzeni
Banacha powiążemy w dalszej części skryptu z zasadą
jednostajnej ograniczoności.
\end{przyklad}

\begin{uwaga}[Dalsza perspektywa: przestrzenie dystrybucji]
W teorii dystrybucji zbiory polarne otoczeń zera w przestrzeni
funkcji testowych są słabo-* zwarte, natomiast topologia silna
na przestrzeni dystrybucji oznacza zbieżność jednostajną
na zbiorach ograniczonych funkcji testowych. Różnica między
zbiorem ograniczonym, zwartym i równociągłym będzie podstawowa
w opisie topologii $\mathcal D'(\Omega)$.
\end{uwaga}

\subsection{Zastosowanie zwartości: nieskończenie wiele warunków liniowych}

\begin{twierdzenie}[Kryterium przedłużenia z wieloma warunkami]
\label{thm:helly-interpolation}
Niech $E$ będzie przestrzenią unormowaną nad $\K$,
$R\ge0$, $I$ dowolnym zbiorem, $(x_i)_{i\in I}$ rodziną
punktów $E$, a $(b_i)_{i\in I}$ rodziną skalarów.
Istnieje funkcjonał $f\in E'$ spełniający
$\|f\|\le R$ i $f(x_i)=b_i$ dla każdego $i\in I$
wtedy i tylko wtedy, gdy dla dowolnego $m\in\N$,
$i_1,\ldots,i_m\in I$ oraz $c_1,\ldots,c_m\in\K$
zachodzi
\[
 \left|\sum_{j=1}^{m}c_jb_{i_j}\right|
 \le R\left\|\sum_{j=1}^{m}c_jx_{i_j}\right\|.
\]
\end{twierdzenie}
\begin{proof}
Konieczność wynika z liniowości funkcjonału i oszacowania
$|f(x)|\le\|f\|\,\|x\|$.

Udowodnijmy dostateczność w sposób wykorzystujący zwartość.
Niech $B_R'=\{f\in E':\|f\|\le R\}$; jest to słabo-* zwarty
zbiór na mocy Wniosku~\ref{cor:alaoglu-norm}. Dla $i\in I$
połóżmy $C_i=\{f\in B_R':f(x_i)=b_i\}$.
Każdy $C_i$ jest słabo-* domknięty w $B_R'$.
Dla ustalonego skończonego $J\subseteq I$ definiujemy
na $M_J=\operatorname{lin}\{x_j:j\in J\}$ funkcjonał
$g_J(\sum_{j\in J}c_jx_j)=\sum_{j\in J}c_jb_j$.
Warunek z tezy zapewnia, że definicja nie zależy od
przedstawienia argumentu, a ponadto $|g_J(x)|\le R\|x\|$.
Twierdzenie Hahna--Banacha daje przedłużenie
$f_J\in E'$ z $\|f_J\|\le R$; stąd
$f_J\in\bigcap_{j\in J} C_j$.
Rodzina $(C_i)_{i\in I}$ ma więc własność skończonych przecięć.
Zwartość $B_R'$ daje $f\in\bigcap_{i\in I}C_i$,
który spełnia wszystkie żądane warunki.
\end{proof}

\begin{przyklad}[Przedłużanie wartości na wektorach współrzędnych]
Niech $E=\ell^1(\N)$, $x_n=e_n$ i $R>0$.
Dla dowolnej rodziny liczb $b_n$ ograniczonej przez
$|b_n|\le R$ otrzymujemy funkcjonał $f\in E'$
spełniający $f(e_n)=b_n$ i $\|f\|\le R$:
\[
 \left|\sum_{j=1}^m c_jb_{i_j}\right|
 \le R\sum_{n=1}^\infty
       \left|\sum_{j:i_j=n}c_j\right|
 =R\left\|\sum_{j=1}^m c_je_{i_j}\right\|_1.
\]
Jest to zgodne z identyfikacją $(\ell^1)'=\ell^\infty$.
\end{przyklad}
